\documentclass{article}%
\usepackage[T1]{fontenc}%
\usepackage[utf8]{inputenc}%
\usepackage{lmodern}%
\usepackage{textcomp}%
\usepackage{lastpage}%
\usepackage{geometry}%
\geometry{margin=2cm}%
\usepackage{expex}%
\usepackage[T1]{fontenc}%
\usepackage[utf8]{inputenc}%
\usepackage[spanish]{babel}%
%
\title{Análisis Lingüístico: Gramatica Normativa Kaqchikel}%
\author{Generado por el TFM de Elías Carrasco}%
%
\begin{document}%
\normalsize%
\maketitle%
\section{Page 89}%
\label{sec:Page89}%
\subsection{1 / 5 / 7 Tajb’en techlal amb’il ojetaq moqa otaq / Uso de ojetaq u otaq}%
\label{subsec:1/5/7Tajbentechlalambilojetaqmoqaotaq/Usodeojetaquotaq}%

%
Se utiliza para indicar que una acción se había terminado antes %
de la ocurrencia de otra acción pero antes del día en que se %
está hablando. %
Ejemplo %
\subsection{1 / 6 Yeky’il tumel / Movimiento y dirección}%
\label{subsec:1/6Yekyiltumel/Movimientoydireccin}%

%
Tienen el papel de indicar la dirección y trayectoria de deter{-} %
minada acción que le describe el hablante al oyente. Todos los %
direccionales se derivan de los llamados verbos intransitivos de %
movimiento. Los direccionales pierden parte de sus elementos %
y son objeto de cambios en su estructura fonológica cuando %
se combinan con otros direccionales o con verbos transitivos. %
(Variación dialectal Mam 2000:110) 

%
\end{document}